\chapter{Summary of Design Decisions}
\label{kap:decisions}

This chapter summarises the key design decisions made during the development of the UACS
board and provides the technical reasoning behind each choice. The goal is to make the
design intent transparent so that future maintainers can understand not only \emph{what}
was chosen, but \emph{why}.

\begin{longtable}{p{0.35\textwidth} p{0.57\textwidth}}
  \caption{Design decisions and their rationale.}
  \label{tab:decisions} \\
  \toprule
  \textbf{Decision} & \textbf{Reasoning} \\
  \midrule
  \endfirsthead

  \multicolumn{2}{c}{\textit{(Table~\ref{tab:decisions} continued)}} \\
  \toprule
  \textbf{Decision} & \textbf{Reasoning} \\
  \midrule
  \endhead

  \bottomrule
  \endlastfoot

  STM32H523CET6 (LQFP-48)
    & Cortex-M33 at \SI{250}{\mega\hertz}; central controller with
      all required peripherals (SPI, I\textsuperscript{2}S, FDCAN, USB)
      in a compact 48-pin package~\cite{ref1}. \\[4pt]

  SWD debug interface (not JTAG)
    & Two-signal interface (SWDIO, SWCLK) provides equivalent debug
      capability while freeing TDI, TDO, and nTRST pins for other
      functions --- a meaningful saving on a pin-limited
      LQFP-48~\cite{ref12,ref13}. \\[4pt]

  SPI for SD card (not SDMMC)
    & The SDMMC peripheral requires pins that are not bonded out on the
      LQFP-48 package; SPI is the only physically available
      option~\cite{ref1}. Package constraint, not design preference. \\[4pt]

  RC522 RFID module
    & Carried over from the predecessor design; existing firmware and
      proven compatibility reduce integration risk~\cite{ref3}. \\[4pt]

  MAX98357A audio amplifier
    & Integrates I\textsuperscript{2}S receiver, Class~D amplifier, and
      direct speaker drive in one device, eliminating a separate DAC,
      amplifier, and output filter. Hardware-selectable gain via
      GAIN\_SLOT~\cite{ref6}. \\[4pt]

  \SI{8}{\ohm} speaker (not \SI{4}{\ohm})
    & Halves the output current relative to a \SI{4}{\ohm} load, keeping
      peak current within the \SI{1}{\ampere} budget of the single buck
      regulator that supplies the entire board. \\[4pt]

  MCP16301H buck converter
    & Carried over from the predecessor design; UACS follows the
      datasheet \SI{3.3}{\volt} reference design closely for a
      known-good result~\cite{ref2}. \\[4pt]

  Dual power input (USB-C + RJ45)
    & RJ45 powers the board during normal deployment from the CAN
      infrastructure; USB-C powers it during development and firmware
      flashing without requiring the full CAN harness. Diode OR-ing
      allows simultaneous connection. \\[4pt]

  TCAN1057A CAN~FD transceiver
    & Meets parameters including CAN~FD support; higher data rate
      (up to 5\,Mbit/s data phase) and 64-byte payload compared with
      classical CAN~\cite{ref14,ref15}. \\[4pt]

  Split CAN termination
    & Maintains \SI{120}{\ohm} differential impedance while adding a
      midpoint capacitor that provides a low-impedance path to ground
      for common-mode noise, improving EMC on a board with a switching
      regulator and long cable run~\cite{ref16,ref17,ref18}. \\[4pt]

  Two \SI{120}{\ohm} resistors in parallel per CAN side
    & \SI{60}{\ohm} resistors in 0805 are costly and less commonly
      stocked; \SI{120}{\ohm} is a standard value and two in parallel
      give exactly \SI{60}{\ohm}. Power dissipation is also split
      between two components. \\[4pt]

  C0G capacitors on CAN bus lines and crystal circuit
    & C0G (Class~1) capacitance is stable with voltage and temperature
      ($\pm30$\,ppm/°C), unlike Class~2 (X7R) which varies with both.
      Stability is critical for oscillator load and for preserving fast
      CAN~FD signal edges~\cite{ref19,ref20}. \\[4pt]

  Ferrite-bead VDDA filter
    & Isolates the MCU analog supply from digital switching noise,
      protecting the ADC reference and internal oscillator. A ferrite
      bead + decoupling capacitor network forms a low-pass filter that
      dissipates high-frequency noise as heat~\cite{ref21}. \\[4pt]

  USB~2.0 device only (no USB~3.x)
    & Carried over from the predecessor design; USB~2.0 full-speed is
      sufficient for DFU firmware flashing and debugging. No higher
      transfer rate is required in this application. \\[4pt]

  External \SI{16}{\mega\hertz} crystal
    & USB full-speed and CAN~FD require a frequency accuracy that the
      internal \SI{64}{\mega\hertz} HSI RC oscillator cannot
      guarantee. An external crystal provides the necessary
      stability~\cite{ref1}. \\[4pt]

  VCAP \SI{2.2}{\micro\farad} capacitors
    & Mandated by the STM32H5 datasheet for stability of the internal
      LDO voltage regulator (\SI{3.3}{\volt}~$\rightarrow$~\SI{1.2}{\volt}
      core); placement and value are part of the regulator's stability
      requirement~\cite{ref1}. \\[4pt]

  LED resistor values from formula
    & $R = (V_\mathrm{DD} - V_f) / I_f$, nearest standard E-series
      value selected, traceable to forward-voltage data in the Kingbright
      datasheet~\cite{ref5}. \\

\end{longtable}
